% Physical unit 035; accepted Session 03 after fresh 400-dpi authority replay.
\GordanUnitStart{035}{printed page 30}
\GordanTopFolio{30}
\begingroup\small
\[
\begin{aligned}
&={1\over\sin A}\cdot{1-2q\over1+q^2}\cdot
 {1-3q^2+4q^2\cos^2{iv\pi\over2K}
  \over 1+2q-4q\cos^2{iv\pi\over2K}}\\
&={1\over\sin A}\cdot{1-2q\over1+q^2}\cdot
 {1-3q^2+q\varepsilon\over 1+2q-\varepsilon{1-4q^2\over1-2q\varepsilon}}\\
&={\sqrt{x'}\over\sin A}\cdot{1-2q\varepsilon\over1+q^2}\cdot {1+q\varepsilon-3q^2\over1-\varepsilon}
\end{aligned}
\]
vel cum sit:
\[
 Civ={1\over C(v,x')}={1\over \sin l_0}
\]
\[
 \sin A={\sqrt{x'}\sin l_0\over1-\varepsilon}\{1-q\varepsilon-4q^2\}
 =\sqrt{\sin l_0\sin B_0}\left\{1+{\varepsilon^2\over2}-q\varepsilon-4q^2\right\}.
\]

\sect{9}

Jam vero alteram formulam computandae longitudinis periodi constituere placet.
Est enim:
\[
\begin{aligned}
1-{2\over\pi}\mathcal A(l_0)&={2K\over\pi}iZ(iv+K)\qquad \text{cap. VI \S\ 4.}\\
&={-4iq\sin {iv\pi\over K}\over 1-2q\cos {iv\pi\over K}}
=-{8iq\tg {iv\pi\over2K}\cos^2{iv\pi\over2K}
\over 1+2q-4q\cos^2{iv\pi\over2K}}\\
&={8q\cos A\over1+2q\sin^2A-2q}
={2\varepsilon\cos A\over {1-2q\varepsilon\over1-2q}-\varepsilon}
=2\varepsilon{1-2q\over1-\varepsilon}\cos A.
\end{aligned}
\]
\[
 \mathcal A(l_0)={\pi\over2}-\pi\varepsilon{1-2q\over1-\varepsilon}\cos A.
\]

\rule{0.35\textwidth}{0.4pt}

\caput{VIII}
\sect{1}

Capite sexto ostendimus, quantitates
\[
 B,\quad l,\quad T,\quad s,\quad \psi,\quad \varphi,
\]
functionibus unius variabilis $u$ determinari posse. Quas ut ad calculos vocemus, $\Theta$ functiones evolvendae sunt in series, progredientes ex ratione et dignitatum constantis $q$ et cosinus atque sinus multiplorum fractionis
\[
 {\pi u\over2K}
\]
quam more analistorum \gq{$x$} significemus.

% Continuation: source PDF pages 036--047.
\endgroup
